\documentclass[UTF8,12pt,a4paper,fontset=fandol]{ctexart}

% 支持从项目根目录或本文件所在目录编译。
\makeatletter
\def\input@path{{../}{../../}}
\makeatother
\usepackage{researchnotes}

\title{行列式自测卷（题目版）}
\author{}
\date{}

\begin{document}
\maketitle

\section*{作答说明}
本卷面向高等代数行列式章节的阶段复习，共 30 题，满分 100 分。
第一部分每题 2 分，第二部分每题 3 分，第三部分每题 5 分。
建议用时 180--240 分钟，也可分三次完成。计算题应写出主要变换，证明题应交代关键依据。
除题目另有说明外，所有矩阵及参数均在实数域 $\mathbb R$ 上，$n$ 为正整数。

记 $I_n$ 为 $n$ 阶单位矩阵，$A^{\mathsf T}$ 为 $A$ 的转置。
对方阵 $A=(a_{ij})$，记 $M_{ij}$ 为删去第 $i$ 行、第 $j$ 列所得的
\keyterm{余子式}（minor），$C_{ij}=(-1)^{i+j}M_{ij}$ 为对应的
\keyterm{代数余子式}（cofactor）。$R_i$、$K_j$ 分别表示第 $i$ 行、第 $j$ 列。
空乘积约定为 $1$，零次幂表示常数多项式 $1$ 的取值。
可使用行列式基本性质、乘法公式、拉普拉斯展开和范德蒙公式；最后两题还需用到一元多项式的因式定理。


\section{基础概念与性质（共 20 分）}
\begin{enumerate}[label=\arabic*.,series=detquestions]

\item\label{q:inversions}
求排列 $(3,6,1,5,2,4)$ 的\keyterm{逆序数}（inversion number），并判断其奇偶性。

\item\label{q:term-sign}
在一般五阶行列式 $\det(a_{ij})_{1\leq i,j\leq5}$ 的完全展开式中，
单项式 $a_{14}a_{21}a_{35}a_{42}a_{53}$ 的系数是多少？说明依据。

\item\label{q:count-terms}
将一般 $n$ 阶行列式看作关于彼此独立的元素 $a_{ij}$ 的多项式，设 $n\geq2$。
完全展开式中含 $a_{12}$ 的项有多少个？同时含 $a_{12}$ 和 $a_{21}$ 的项有多少个？

\item\label{q:row-operations}
设三阶矩阵 $A$ 满足 $\det A=5$。从 $A$ 出发，依次交换第一、三行，
将第二行乘以 $-2$，再把第一行的 $3$ 倍加到第三行，所得矩阵记为 $B$。求 $\det B$。

\item\label{q:minor-cofactor}
设
\[
A=\begin{pmatrix}1&2&0\\3&-1&4\\2&5&1\end{pmatrix}.
\]
求 $M_{23}$ 和 $C_{23}$，并写出二者的符号关系。

\item\label{q:cofactor-sums}
设三阶矩阵 $A=(a_{ij})$ 满足 $\det A=7$。求
\[
\sum_{j=1}^{3}a_{2j}C_{2j},\qquad
\sum_{j=1}^{3}a_{1j}C_{2j},
\]
并简述理由。

\item\label{q:triangular}
计算下列两个行列式，并说明第二个行列式的符号：
\[
D_1=\begin{vmatrix}-2&1&0&3\\0&3&4&1\\0&0&1&5\\0&0&0&4\end{vmatrix},
\qquad
D_2=\begin{vmatrix}0&0&0&1\\0&0&2&0\\0&3&0&0\\4&0&0&0\end{vmatrix}.
\]

\item\label{q:dependent-rows}
设 $D_n=\det(i+j)_{1\leq i,j\leq n}$，$n\geq2$。分别求 $n=2$ 与 $n\geq3$ 时的 $D_n$。

\item\label{q:true-false}
判断下列断言的正误；正确的说明理由，错误的给出反例。
\begin{enumerate}[label=(\arabic*)]
\item 对任意同阶方阵 $A,B$，都有 $\det(A+B)=\det A+\det B$。
\item 一个三阶行列式等于零，则它一定有两行成比例。
\item 对任意实方阵 $A$，都有 $\det(A^{\mathsf T}A)\geq0$。
\end{enumerate}

\item\label{q:product-scaling}
设 $A,B$ 均为三阶矩阵，$\det A=-2$，$\det B=3$。求
\[
\det(2A),\qquad \det(A^{\mathsf T}B),\qquad \det(A^2).
\]
\end{enumerate}

\section{计算与方法训练（共 30 分）}
\begin{enumerate}[label=\arabic*.,resume=detquestions]

\item\label{q:numeric-three}
用初等行变换计算
\[
D=\begin{vmatrix}1&2&3\\2&5&7\\1&0&4\end{vmatrix}.
\]

\item\label{q:numeric-four}
计算
\[
D=\begin{vmatrix}
1&1&1&1\\1&2&1&1\\1&1&3&1\\1&1&1&4
\end{vmatrix}.
\]

\item\label{q:parameter-three}
将下列行列式分解为一次因式的乘积，并求它为零时 $x$ 的全部实数取值：
\[
D(x)=\begin{vmatrix}x&1&1\\1&x&1\\1&1&x\end{vmatrix}.
\]

\item\label{q:vandermonde-numeric}
利用\keyterm{范德蒙行列式}（Vandermonde determinant）计算
\[
D=\begin{vmatrix}
1&1&1&1\\-1&0&2&3\\1&0&4&9\\-1&0&8&27
\end{vmatrix}.
\]

\item\label{q:polynomial-four}
设 $x_1,x_2,x_3,x_4\in\mathbb R$。将下列行列式写成乘积形式：
\[
D=\begin{vmatrix}
1&1&1&1\\
x_1+1&x_2+1&x_3+1&x_4+1\\
x_1^2-x_1+2&x_2^2-x_2+2&x_3^2-x_3+2&x_4^2-x_4+2\\
x_1^3+2x_1&x_2^3+2x_2&x_3^3+2x_3&x_4^3+2x_4
\end{vmatrix}.
\]

\item\label{q:constant-off-diagonal}
设 $n\geq2$，$a,b\in\mathbb R$，矩阵 $A_n=(a_{ij})$ 的元素为
\[
a_{ij}=\begin{cases}a,&i=j,\\b,&i\ne j.\end{cases}
\]
求 $\det A_n$，并说明所得公式在 $a=b$ 时是否成立。

\item\label{q:arrow}
设 $n\geq2$，计算下列 $n$ 阶行列式，其中 $d_2,\ldots,d_n$ 允许为零：
\[
D_n=\begin{vmatrix}
\alpha&\beta_2&\beta_3&\cdots&\beta_n\\
\gamma_2&d_2&0&\cdots&0\\
\gamma_3&0&d_3&\cdots&0\\
\vdots&\vdots&\vdots&\ddots&\vdots\\
\gamma_n&0&0&\cdots&d_n
\end{vmatrix}.
\]

\item\label{q:tridiagonal}
设 $T_n=(t_{ij})_{1\leq i,j\leq n}$，其中
\[
t_{ij}=\begin{cases}2,&i=j,\\-1,&|i-j|=1,\\0,&|i-j|\geq2.\end{cases}
\]
记 $D_n=\det T_n$。建立递推关系并求 $D_n$ 的通项。

\item\label{q:cyclic}
设 $n\geq3$，矩阵 $A=(a_{ij})$ 仅在以下位置可能非零：
\[
a_{ii}=u_i\ (1\leq i\leq n),\qquad
a_{i,i+1}=v_i\ (1\leq i<n),\qquad a_{n1}=v_n.
\]
其余元素均为零。求 $\det A$。

\item\label{q:laplace}
使用\keyterm{拉普拉斯展开}（Laplace expansion），按第一、三行展开并计算
\[
D=\begin{vmatrix}
1&2&0&0\\0&0&3&4\\5&6&0&0\\0&0&7&8
\end{vmatrix}.
\]
请写出所选子式及其代数余子式的符号。
\end{enumerate}

\section{综合应用与证明（共 50 分）}
\begin{enumerate}[label=\arabic*.,resume=detquestions]

\item\label{q:cramer}
用\keyterm{克拉默法则}（Cramer's rule）求解方程组，写出系数行列式及三个替换列行列式：
\[
\begin{cases}
x+y+z=6,\\
2x-y+3z=9,\\
3x+2y-z=4.
\end{cases}
\]

\item\label{q:parameter-system}
讨论实参数 $\lambda$ 取不同值时，下列方程组的解的情况；有解时写出全部解：
\[
\begin{cases}
\lambda x+y+z=1,\\
x+\lambda y+z=1,\\
x+y+\lambda z=-2.
\end{cases}
\]

\item\label{q:all-cofactors}
设 $n\geq2$，$J_n$ 为所有元素都等于 $1$ 的 $n$ 阶矩阵。
对任意 $n$ 阶矩阵 $A$，证明
\[
\det(A+tJ_n)=\det A+t\sum_{i=1}^n\sum_{j=1}^n C_{ij}
\qquad(t\in\mathbb R),
\]
其中 $C_{ij}$ 始终指原矩阵 $A$ 的代数余子式。进而求
\[
A=\begin{pmatrix}2&1&0\\0&3&1\\1&0&2\end{pmatrix}
\]
的全部代数余子式之和，要求不逐个计算九个代数余子式。

\item\label{q:skew-symmetric}
设实方阵 $A$ 满足 $A^{\mathsf T}=-A$，称其为
\keyterm{反对称矩阵}（skew-symmetric matrix）。
证明：若 $A$ 的阶数为奇数，则 $\det A=0$。
再给出一个偶数阶反对称矩阵，其行列式不为零。

\item\label{q:block-determinant}
设 $A,B$ 为任意两个 $n$ 阶实矩阵。用行、列变换证明
\[
\det\begin{pmatrix}A&B\\B&A\end{pmatrix}
=\det(A+B)\det(A-B).
\]
不假设 $A$、$B$ 可逆，也不假设 $AB=BA$。

\item\label{q:rank-one}
设 $u=(u_1,\ldots,u_n)^{\mathsf T}$、$v=(v_1,\ldots,v_n)^{\mathsf T}$ 为实列向量，
$uv^{\mathsf T}$ 表示第 $(i,j)$ 个元素为 $u_iv_j$ 的矩阵。
用行列式的多重线性性质证明
\[
\det(I_n+uv^{\mathsf T})=1+\sum_{i=1}^n u_iv_i.
\]
进而说明，对任意实向量 $u$，$\det(I_n+uu^{\mathsf T})>0$。

\item\label{q:min-index}
设 $s_1,\ldots,s_n\in\mathbb R$，$A=(s_{\min(i,j)})_{1\leq i,j\leq n}$。
求 $\det A$，并给出它不为零的充要条件。
例如 $n=4$ 时，
\[
A=\begin{pmatrix}
s_1&s_1&s_1&s_1\\
s_1&s_2&s_2&s_2\\
s_1&s_2&s_3&s_3\\
s_1&s_2&s_3&s_4
\end{pmatrix}.
\]

\item\label{q:polynomial-basis}
设 $p_k(t)=\sum_{r=0}^{k}c_{kr}t^r$（$k=0,1,\ldots,n-1$）为实系数多项式，
允许 $c_{kk}=0$。对任意 $x_1,\ldots,x_n\in\mathbb R$，证明
\[
\det\bigl(p_{i-1}(x_j)\bigr)_{1\leq i,j\leq n}
=\left(\prod_{k=0}^{n-1}c_{kk}\right)
\prod_{1\leq i<j\leq n}(x_j-x_i).
\]
进而求
\[
\begin{vmatrix}
2&2&2\\
3a-1&3b-1&3c-1\\
5a^2+a+4&5b^2+b+4&5c^2+c+4
\end{vmatrix}.
\]

\item\label{q:missing-power}
设 $n\geq2$，$x_1,\ldots,x_n\in\mathbb R$。
矩阵 $W$ 的前 $n-1$ 行依次为 $(x_j^k)_{j=1}^n$（$k=0,\ldots,n-2$），
最后一行为 $(x_j^n)_{j=1}^n$，即跳过了幂次 $n-1$。
证明
\[
\det W=\left(\sum_{j=1}^n x_j\right)
\prod_{1\leq i<j\leq n}(x_j-x_i).
\]
例如 $n=2$ 时，$W=\begin{pmatrix}1&1\\x_1^2&x_2^2\end{pmatrix}$。

\item\label{q:interpolation-sum}
设 $n\geq2$，$x_1,\ldots,x_n$ 是两两不同的实数。
证明：对每个整数 $m$，$0\leq m\leq n-1$，都有
\[
\sum_{k=1}^n
\frac{x_k^m}{\displaystyle\prod_{\substack{1\leq j\leq n\\j\ne k}}(x_k-x_j)}
=\begin{cases}0,&0\leq m\leq n-2,\\1,&m=n-1.\end{cases}
\]
要求用范德蒙行列式及代数余子式展开证明，并说明为什么需要节点两两不同。
\end{enumerate}

\end{document}
